% !TEX root = ../algebra_02.tex

\section{Ранг произведения матриц. Связь ранга с \\ PDQ-разложением}

\subsection{Ранг произведения матриц}

\begin{Theorem}{Неравенства для рангов}{}
	Для матриц $A \in k^{m \times n}$, $B \in k^{n \times p}$:
	\begin{enumerate}
		\item $\mathrm{rk}(A + B) \leq \mathrm{rk}(A) + \mathrm{rk}(B)$.
		\item $\mathrm{rk}(AB) \leq \min(\mathrm{rk}(A),\, \mathrm{rk}(B))$.
		\item Если $U \in GL_m(k)$, то $\mathrm{rk}(UA) = \mathrm{rk}(A)$. Если $V \in GL_n(k)$, то $\mathrm{rk}(AV) = \mathrm{rk}(A)$.
	\end{enumerate}
\end{Theorem}

\begin{proof}
	\begin{enumerate}
		\item Пункт 1 следует из $\mathrm{Lin}(A+B) \subseteq \mathrm{Lin}(A) + \mathrm{Lin}(B)$.
		\item По определению матричного умножения, каждый столбец матрицы $AB$ является линейной комбинацией столбцов матрицы $A$. Поэтому столбцы $AB$ лежат в линейной оболочке столбцов $A$, откуда $\mathrm{rk}(AB) \leq \mathrm{rk}(A)$.
		Оценка $\mathrm{rk}(AB) \leq \mathrm{rk}(B)$ следует из операции транспонирования: 
		\[
		\mathrm{rk}(AB) = \mathrm{rk}((AB)^T) = \mathrm{rk}(B^T A^T) \leq \mathrm{rk}(B^T) = \mathrm{rk}(B).
		\]
		Объединяя два неравенства, получаем $\mathrm{rk}(AB) \leq \min(\mathrm{rk}(A),\, \mathrm{rk}(B))$.
		
		\item Это следствие пункта 2: $\mathrm{rk}(UA) \leq \mathrm{rk}(A)$. С другой стороны, так как $U$ обратима, мы можем записать $A = U^{-1}(UA)$, откуда $\mathrm{rk}(A) \leq \mathrm{rk}(UA)$. Следовательно, $\mathrm{rk}(UA) = \mathrm{rk}(A)$.
		Аналогично, $\mathrm{rk}(AV) \leq \mathrm{rk}(A)$ и $\mathrm{rk}(A) = \mathrm{rk}((AV)V^{-1}) \leq \mathrm{rk}(AV)$, поэтому $\mathrm{rk}(AV) = \mathrm{rk}(A)$.
	\end{enumerate}
\end{proof}

\subsection{Связь ранга с PDQ-разложением}

\begin{Theorem}{PDQ-разложение и ранг}{}
	Любая матрица $A \in k^{m \times n}$ ранга $r$ допускает представление в виде
	\[
	A = P D Q,
	\]
	где $P \in GL_m(k)$ и $Q \in GL_n(k)$ — обратимые матрицы, а матрица $D \in k^{m \times n}$ имеет нормальный блочный вид
	\[
	D = \begin{pmatrix} E_r & 0 \\ 0 & 0 \end{pmatrix},
	\]
	где $E_r$ — единичная матрица порядка $r$. Ранг матрицы $A$ в точности равен размеру блока $E_r$.
\end{Theorem}

\begin{proof}
	Матрицу $A$ можно элементарными преобразованиями строк и столбцов привести к окаймлённой единичной матрице $D = \begin{pmatrix} E_r & 0 \\ 0 & 0 \end{pmatrix}$. 
	Каждое элементарное преобразование строк равносильно умножению слева на некоторую обратимую матрицу, а преобразование столбцов — умножению справа. Следовательно, существуют обратимые матрицы $U \in GL_m(k)$ и $V \in GL_n(k)$ такие, что $U A V = D$. 
	
	Отсюда выражается сама матрица: $A = U^{-1} D V^{-1}$. Полагая $P = U^{-1}$ и $Q = V^{-1}$, получаем искомое разложение $A = P D Q$.
	Поскольку умножение на обратимые матрицы не меняет ранг, $\mathrm{rk}(A) = \mathrm{rk}(D)$. Ранг матрицы $D$ равен количеству единиц на главной диагонали, то есть $r$. Таким образом, ранг неразрывно связан с размером единичного блока в PDQ-разложении.
\end{proof}
